\uuid{nQsF}
\exo7id{3318}
\titre{exo7 3318}
\auteur{quercia}
\organisation{exo7}
\datecreate{2010-03-09}
\isIndication{false}
\isCorrection{true}
\chapitre{Espace vectoriel}
\sousChapitre{Base}
\cpge{mpsi-mp2i.espaces-vectoriels.c, pcsi.espaces-vectoriels.c, ptsi.espaces-vectoriels.c, mpsi-mp2i.changements-bases.b, mpsi-mp2i.matrices-applications-lineaires.c, pcsi.matrices-applications-lineaires.d, ptsi.matrices-applications-lineaires.d, bcpst1.espace-vectoriel-kn.b, bcpst1.systemes-lineaires.c}
\module{Algèbre}
\niveau{L1}
\difficulte{}

\contenu{
\question{
Dans $\R^4$, trouver le rang de la famille de vecteurs :

$$\vec a = (3,2,1,0),\quad
  \vec b = (2,3,4,5),\quad
  \vec c = (0,1,2,3),\quad
  \vec d = (1,2,1,2),\quad
  \vec e = (0,-1,2,1).$$
}
\reponse{
$r = 3$,\quad $2\vec a - 3\vec b + 5\vec c = \vec 0$,
                 \quad $\vec b - 2\vec d - \vec e = \vec 0$.
}
}
